% =====================================================================
%  Phase C --- Route optimisation and the energy budget
% =====================================================================
\section{Route Optimisation and the Energy Budget}
\label{sec:phaseC_route}

\subsection{The cost of the obvious order}

The Cloud lists stations in catalogue order --- \station{P1,1R},
\station{P1,1L}, \station{P1,2R}, \ldots{} --- which is right for a human
reading a list and wrong for a robot driving a building. The two sides of
one plant are in \emph{different bands} (Table~\ref{tab:bands}), separated
by an \SI{8}{m} gutter, so \station{R} to \station{L} means driving to the
end of the row, crossing at a headland, and coming back. Over 48 stations
the catalogue order asks for that 45 times.

Measured against the same \texttt{route()} the driver actually flies,
starting from the dock:

\begin{table}[!t]
\caption{Route cost over all 48 stations. Distances are computed
along the real leg geometry, not estimated.}
\label{tab:route}
\centering\footnotesize
\begin{tabular}{@{}lrr@{}}
\toprule
order & distance & band crossings \\
\midrule
catalogue order, as issued & \measured{312~m} & \measured{45} \\
swept band by band         & \measured{40~m}  & \measured{3}  \\
\midrule
saving                     & \measured{272~m} (\measured{87\%}) & 42 \\
\bottomrule
\end{tabular}
\end{table}

The signature of the defect was visible in the campaign log before it was
diagnosed. Inter-station intervals within a row rose and then fell ---
$23, 28, 31, 34, 38, 39, 35, \ldots, 27, 24$~s --- a bell curve, because
a plant near a headland is cheap to cross at and a plant in the middle of
the row is expensive from either end.

\subsection{The order chosen, and why it is not a solver}

Within a band the robot drives freely along $x$; leaving one costs a
headland trip whatever the destination. The cheap route is therefore the
obvious one --- sweep a band end to end, cross once, sweep the next back
--- and the only remaining decisions are which end to start at and which
direction to take. That is four combinations. All four are evaluated
exactly, against the real route geometry, and the shortest wins.

A TSP solver would spend milliseconds rediscovering that the optimal tour
of points on a line is to walk the line, and would then be free to return
something else the day a constant changed.

\paragraph{The property that makes it safe to enable} \textbf{The order
as issued is the fifth candidate, and it wins ties.} Adopting a longer
route is therefore not a bug that could occur --- it is something the
function cannot do. The suite checks it over \measured{720} random
subsets of every size, and separately asserts that a request too small to
benefit is returned untouched and that an unknown label does not raise.

\paragraph{The objection that was answered rather than accepted} The
previous implementation carried a note saying reordering was omitted so
that the acknowledgement stream would keep matching the request. That
was a real concern and it is addressed directly: the mission retains the
issued order alongside the driven one, every acknowledgement still names
its own station, and the Cloud tracks progress by label. Nothing an
operator reads was ever indexed by position.

\subsection{Measured effect}

\begin{table}[!t]
\caption{Full 48-station campaign, before and after reordering. Times
are wall-clock, from the Cloud's own request record.}
\label{tab:campaign}
\centering\footnotesize
\begin{tabular}{@{}lrr@{}}
\toprule
 & catalogue order & swept order \\
\midrule
campaign duration     & \measured{1\,505~s} & \measured{601~s} \\
stations measured     & \measured{48/48}    & \measured{48/48} \\
reports rejected      & \measured{0}        & \measured{0}     \\
floor camera resolved & \measured{48/48}    & \measured{48/48} \\
\bottomrule
\end{tabular}
\end{table}

The campaign runs in \measured{40\%} of the time for identical coverage
and identical data quality.

\subsection{Distance is measured; energy is modelled}

This distinction is maintained in the code, in the CSV column names, and
here, because collapsing it would put a fabricated number beside
forty-eight real ones.

\paragraph{Measured} The distance driven is integrated from odometry in
the driver, summed over successive positions rather than from the
reported speed --- the speed is sampled at 14~Hz and each sample would be
held for the whole interval, accumulating a different error than the path
itself has. A single step longer than \SI{1.0}{m} is a respawn, not a
drive, and is not counted. The figure inherits the odometry's drift, which
is the honest behaviour: it measures what the robot \emph{believes} it
drove, which is the right basis for an energy estimate built on it.

\paragraph{Modelled} Gazebo models no battery and the platform carries no
current sensor. Rather than display a percentage ticking down,
\filename{agri/survey.py} states four constants in the open, labelled in
the source as \textsc{assumptions, not measurements}:

\begin{equation}
E = P_{\text{idle}} \cdot t_{\text{mission}}
  + P_{\text{drive}} \cdot t_{\text{driving}},
\qquad
t_{\text{driving}} = \frac{d}{v_{\text{cruise}}}
\label{eq:energy}
\end{equation}

with $P_{\text{idle}} = \SI{18}{W}$ (computer, lidar, two cameras, MQTT
link), $P_{\text{drive}} = \SI{45}{W}$ (four wheel motors, while moving),
$v_{\text{cruise}} = \SI{0.45}{m/s}$ and a nameplate pack of
\SI{120}{Wh}. The suite asserts $v_{\text{cruise}}$ equals the driver's
own \param{V\_MAX}, so the estimate cannot drift away from the robot it
describes, and asserts that $t_{\text{driving}}$ can never exceed the
mission duration --- otherwise a bad odometry reading would invent drive
energy out of nothing.

\subsection{The finding that matters more than the 87\%}

Applying Eq.~\ref{eq:energy} to the two campaigns:

\begin{table}[!t]
\caption{Energy budget. Distance and duration are measured; the three
energy columns are derived from Eq.~\ref{eq:energy}.}
\label{tab:energy}
\centering\footnotesize
\begin{tabular}{@{}lrrrr@{}}
\toprule
campaign & $d$ & $t$ & idle & drive \\
\midrule
catalogue order & \measured{312~m} & \measured{1505~s} & \modelled{7.5~Wh} & \modelled{8.7~Wh} \\
swept order     & \measured{40~m}  & \measured{601~s}  & \modelled{3.0~Wh} & \modelled{1.1~Wh} \\
\bottomrule
\end{tabular}
\\[2pt]\modelnote
\end{table}

Idle power is drawn whether or not the robot moves. On a 48-station
sweep the robot spends far longer standing at marks --- settling,
focusing, photographing, encoding a QR code, sealing an envelope --- than
driving between them, so \textbf{shortening the route removes drive
energy and leaves idle energy untouched: it attacks the smaller half of
the bill.}

This is the more useful conclusion for a deployment. Having taken the
easy 87\%, further reduction is no longer a path-planning problem at all;
it is a question of per-station dwell time, and therefore of the vision
settling tolerance, the number of correction attempts, and the image
encoding cost. Sec.~\ref{sec:discussion} returns to it.

\subsection{What is exported}

\filename{requests.csv} keeps measured and modelled apart by name:
\param{driven\_m} and \param{planned\_m} beside \param{energy\_wh\_est},
\param{idle\_wh\_est} and \param{drive\_wh\_est}. A request whose robot
reported no distance leaves those cells \textbf{empty}, not zero --- ``this
robot does not count metres'' and ``it drove nothing'' are different
facts, and the suite asserts the empty cells directly.
